\documentclass{article}

\usepackage{array}
\usepackage{tabularx}
\usepackage[roman,breakwithin,reset]{parnotes}
\usepackage{supertabular}
\usepackage{perpage}
\usepackage{lipsum}

\title{\Huge Building Tables Part V}
\date{}

\MakePerPage{footnote} % May take two runs to work
\setlength\extrarowheight{3pt}

\begin{document}
	\maketitle
	
	\section{Footnotes in the tabular Environment}
	
	\subsection{First Case: using footnote mark and 			text commands}
	\begin{center}
		\begin{tabular}{|c|c|}
			\hline A & B \footnotemark \\
			\hline C & D \\ \hline
		\end{tabular}
		\footnotetext{For the standalone tabular
		env. only one footnote is allowed.}	
	\end{center}	
	
	\bigskip
	
	\subsection{Second Case: using footnote cmd with 			minipage}
	\begin{minipage}{\linewidth}
		\hspace{5cm}
		\begin{tabular}{|c|c|}
			\hline A \footnote{Multiple footnote
				commands are allowed for a tabular
				inside a minipage.} & B \\
			\hline C & D  \footnote{Another footnote.} 				\\ \hline
		\end{tabular}
	\end{minipage}
	
	\bigskip	
	
	\subsection{Third Case: when wrapped in a float}
	\begin{table}[h]
		\centering
		\begin{tabular}{|c|c|}
			\hline A & B \footnotemark[2] \\
			\hline C & D \footnotemark[3] \\ \hline
		\end{tabular}	
	\end{table}
	\footnotetext[2]{Footnote text command must be 				outside the float environment.}
	\footnotetext[3]{Second footnote in the same table.}
	
	\newpage % \pagebreak doesn't add space
	
	\section{Footnotes in the tabularx Environment}
	
	\subsection{footnote cmd in Standalone tabularx}
	\hspace{2.9cm}
	\begin{tabularx}{0.5\linewidth}{|X|X|X|}
		\hline
		{\bf Item} & {\bf Price} & {\bf Quantity}\\
		\hline
		No.1 & 150 & 4 
			\footnote{tabularx Standalone Footnote} \\
		No.2 & 130 & 5 \\
		No.3 & 140 & 6 \\
		\hline
	\end{tabularx} 
	
	\subsection{tabularx Inside table Float}
	\begin{table}[h]
		\centering
		\begin{tabularx}{0.5\linewidth}{|X|X|X|}
		\hline
		{\bf Item} & {\bf Price} & {\bf Quantity}\\
		\hline
		No.1 & 150 & 4 \footnotemark \\
		No.2 & 130 & 5 \\
		No.3 & 140 & 6 \\
		\hline
		\end{tabularx}
	\end{table}
	\footnotetext{This is a footnote for 
			a tabularx inside table float.}
			
	\subsection{tabularx and parnotes Package}
	\begin{table}[h] 
		\centering 
		\begin{tabularx}{0.5\linewidth}{|X|X|X|}
		\hline
		{\bf Item} & {\bf Price} & {\bf Quantity}\\
		\hline
		No.1 & 150 & 4 
					\parnote{This is a parnote!} \\
		No.2 & 130 & 5 
					\parnote{Yet another parnote!}\\
		No.3 & 140 & 6 \\
		\hline
		\end{tabularx}
		\parnotes 
	\end{table}
	
	\newpage
	
	\section{Footnotes for Multipage Tables}
	
	\subsection{Inside longtable} 
	You can directly use the \emph{footnote} command
		after any entry in this environment.
		
	\subsection{For the supertabulars}
	\tablecaption{SuperTabular Above Caption.}
	\tablehead{
		\hline
		\textbf{Counter} & \textbf{Data} &
		\textbf{Data} & \textbf{Data} & \textbf{Data}\\ 
		\hline
		}
	\tabletail{\hline}
	\begin{center}
	\begin{mpsupertabular}{|*{5}{c|}}
		Item & 1 & 2 & 3 & 4 \\
		Item & 1 & 2 & 3 & 4 \\
		Item & 1 & 2 & 3 & 4 \\
		Item & 1 & 2 & 3 & 4 \\
		Item & 1 & 2 & 3 & 4 \\
		Item & 1 & 2 & 3 & 4 \\
		Item & 1 & 2 & 3 & 4 \\
		Item & 1 & 2 & 3 & 4 \\
		Item & 1 & 2 & 3 & 4 \\
		Item & 1 & 2 & 3 & 4 
			\footnote{Supertabular Footnote} \\
		Item & 1 & 2 & 3 & 4 \\
		Item & 1 & 2 & 3 & 4 \\
		Item & 1 & 2 & 3 & 4 \\
		Item & 1 & 2 & 3 & 4 \\
		Item & 1 & 2 & 3 & 4 \\
		Item & 1 & 2 & 3 & 4 \\
		Item & 1 & 2 & 3 & 4 \\
		Item & 1 & 2 & 3 & 4 \\
		Item & 1 & 2 & 3 & 4 \\
		Item & 1 & 2 & 3 & 4 
			\parnote{A parnote inside supertabular} \\
		Item & 1 & 2 & 3 & 4 \\
		Item & 1 & 2 & 3 & 4 \\
		Item & 1 & 2 & 3 & 4 \\
		Item & 1 & 2 & 3 & 4 \\
		Item & 1 & 2 & 3 & 4 \\
		Item & 1 & 2 & 3 & 4 \\
		Item & 1 & 2 & 3 & 4 \\
		Item & 1 & 2 & 3 & 4 \\
		Item & 1 & 2 & 3 & 4 \\
		Item & 1 & 2 & 3 & 4 \\
		Item & 1 & 2 & 3 & 4 \\
		Item & 1 & 2 & 3 & 4 \\
	\end{mpsupertabular}
	\parnotes
	\end{center}
\end{document}
